\section{Mechanism design and auctions}
\label{sec:mechanism-auctions}

The mechanism-design and auction layers are direct consumers of
the kernel-game core.  They serve a dual role:\ they formalize
classical results that belong in any game-theory library, and
they smoke-test the reusability of the core.  Bayesian games and the
auction payoff layers build on the kernel-game/game-form layer and the
solution-concept catalogue; the quasilinear mechanism and
social-choice layers introduce their own finite structures (payments,
contracts, signal structures, bargaining, social-welfare functions)
while reusing the preference and solution-concept vocabulary.

\subsection{Bayesian games}
\label{sec:mech:bayesian}

A \emph{Bayesian game}~\citep{Harsanyi1967} models situations
where players have \emph{private information} that affects payoffs
— each player knows their own private state but only has a prior
belief over others'.  Harsanyi's reduction encodes this private
information as a \emph{type} $\theta_i \in \Theta_i$ for each
player; types are drawn jointly from a common-prior distribution
$\mu \in \PMF{\prod_i \Theta_i}$ known to all players, but only
$\theta_i$ is observed by player $i$.  A pure strategy is then a
contingent plan $\Theta_i \to \Act_i$ that prescribes an action
for every type the player might turn out to be.  Concretely, a
finite Bayesian game is the tuple
$(\Theta, \Act, \mu, u)$ with finite $\Theta_i, \Act_i$ and a
payoff $u : (\prod_i \Theta_i) \times (\prod_i \Act_i) \to
(\Players \to \RR)$ depending on both the realized type vector and
the chosen joint action.  Behavioral strategies replace the pure
version by
$\Theta_i \to \PMF{\Act_i}$.

The Bayesian game compiles to a kernel game via the \TT{ofEU}
adapter (\Cref{sec:semantic-core:kernelgame}), which packages an
EU-only specification as a kernel game whose point-mass outcome
carries the expected-utility payoff vector.  The strategy type is
$\Theta_i \to \Act_i$, and the outcome kernel is the deterministic
point mass at the \emph{ex-ante} expected utility
\[
  \mathrm{exAnteEU}(B, \Profile, i)
    \;=\; \sum_{\theta} \mu(\theta)\;
        u(\theta, \Profile_1(\theta_1), \dots, \Profile_n(\theta_n))(i),
\]
folded into a single payoff vector per profile.  The
type$\,\times\,$action sampling is therefore not a separate
sampling step in the kernel game's $\K$ but is folded into the EU
function — a deliberate use of \TT{ofEU} that lets every
EU-specific solution concept on \KG\ apply to Bayesian games
without further bookkeeping.

\emph{Bayesian Nash equilibrium}~\citep{Harsanyi1967} is then the
ordinary $\IsNash$ predicate of the compiled kernel game,
unfolded as
\[
  \forall i,\, \forall s' : \Theta_i \to \Act_i.\quad
    \mathrm{exAnteEU}(B, \Profile, i) \;\ge\;
    \mathrm{exAnteEU}(B, \Profile[i \mapsto s'], i).
\]
This is the ex-ante Bayes-Nash condition used by the current
library API; ex-post and interim variants are not separate packaged
definitions in this file.

\tradeoff{Continuous type spaces and continuous action spaces are
the natural setting for much of auction theory beyond
sealed-bid mechanisms with discrete bid grids.  As discussed in
\Cref{sec:prelim:pmf}, this falls outside the discrete-PMF
discipline and is not in scope.  All Bayesian-game results in the
library hold for finite type and action spaces.}

\leancorr{\TT{Mechanism/BayesianGame.lean} (\TT{BayesianGame},
\TT{exAnteEU}, \TT{toKernelGame}, \TT{BayesNash},
\TT{bayesNash\_iff\_exAnteEU}).  Supplement~§8.1.}

\subsection{Finite information design}
\label{sec:mech:info-design}

The information-design layer packages the finite Bayesian-persuasion
substrate used by signaling and disclosure examples.  A
\emph{signal structure} is a finite stochastic kernel
$S : \Omega \to \PMF{M}$ from states to public messages.  Given a
prior $\mu \in \PMF{\Omega}$, it induces the joint law
\[
  \nu_S(\omega,m) = \mu(\omega) S(m \mid \omega)
\]
over states and messages, together with the message marginal
$\nu_S^M$, whose finite-sum form is recorded as
\TT{messageMarginal\_apply}.  The predicate \TT{HasPriorMarginal}
says that a joint state-message distribution has state marginal
$\mu$; the basic theorem \TT{joint\_hasPriorMarginal} proves that
every signal structure satisfies this identity.  The library also
exposes two canonical signals:\ \TT{uninformative}, which sends the
unique message regardless of the state, and \TT{fullInformation},
which reports the realized state.

A finite persuasion problem fixes a prior, a signal structure, a
sender payoff, and a receiver payoff.  A decision rule maps public
messages to receiver actions.  Receiver optimality is written with
the unnormalized message weights
\[
  \sum_{\omega \in \Omega}
    \mu(\omega) S(m \mid \omega) u_R(\omega,a),
\]
which is equivalent to posterior expected utility whenever message
$m$ has positive probability and remains meaningful at zero-probability
messages.  A rule is \TT{IsPersuasive} when every induced action is
receiver-optimal at its message.  Sender welfare is the expected
sender utility under the induced joint distribution, with the
finite-sum form recorded as \TT{senderEU\_eq\_sum}; equivalently,
it is the sum of message-contingent sender scores
(\TT{senderEU\_eq\_sum\_senderScore}).

The signal-structure layer is intentionally a substrate rather than the
concavification (optimal-value) theorem:\ it gives the finite objects,
Bayes-plausibility invariant, persuasion predicate, and sender
objective on which finite persuasion examples and later optimal-design
theorems can be stated.

\leancorr{\TT{Mechanism/InformationDesign.lean}
(\TT{SignalStructure}, \TT{joint}, \TT{HasPriorMarginal},
\TT{joint\_hasPriorMarginal}, \TT{messageMarginal\_apply},
\TT{PersuasionProblem}, \TT{IsPersuasive},
\TT{senderEU\_eq\_sum}, \TT{senderEU\_eq\_sum\_senderScore}).}

\paragraph{Feasible posteriors and the splitting characterization.}\
A complementary, belief-based view of disclosure is packaged in
\TT{Mechanism/FeasiblePosteriors.lean}, stated fully generally with no
finiteness hypotheses.  An \emph{experiment} is a distribution
$\tau \in \PMF{\PMF{\Omega}}$ over posterior beliefs; its
\emph{mean posterior} $\overline{\tau} = \tau \mathbin{\bind}
\mathrm{id}$ averages those beliefs back to a prior.  The experiment
is \emph{Bayes-plausible} for $\mu$ exactly when
$\overline{\tau} = \mu$ — the splitting/martingale condition of
\citet{KamenicaGentzkow2011}.  Rather than route through
measure-theoretic conditioning, the file builds the canonical coupling
that realizes a Bayes-plausible experiment (the joint law
$\tau(b)\,b(\omega)$ whose belief marginal is $\tau$ and whose state
marginal is $\overline{\tau}$), proves the splitting direction
(\TT{isBayesPlausible\_iff\_coupling\_fst}), and records that the
feasible set is convex (\TT{isBayesPlausible\_bind}) and closed under
sequential composition (\TT{isBayesPlausible\_bind\_pointwise}), with
the uninformative and fully-revealing experiments as the extremes.

The multi-receiver generalization
(\TT{Mechanism/JointFeasiblePosteriors.lean}) calls a joint experiment
over profiles of posteriors \emph{feasible} when a single common-prior
coupling induces every agent's beliefs at once.  The library proves
the necessary direction — every feasible joint experiment has each
agent's marginal Bayes-plausible
(\TT{IsJointFeasible.bayesPlausible\_agentMarginal}) — and is
deliberately honest that the converse fails for two or more agents
\citep{Arieli2021}, so no full characterization is claimed.

\leancorr{\TT{Mechanism/FeasiblePosteriors.lean}
(\TT{meanPosterior}, \TT{IsBayesPlausible}, \TT{posteriorCoupling},
\TT{isBayesPlausible\_iff\_coupling\_fst},
\TT{isBayesPlausible\_bind});
\TT{Mechanism/JointFeasiblePosteriors.lean}
(\TT{IsJointFeasible},
\TT{IsJointFeasible.bayesPlausible\_agentMarginal}).}

\subsection{Mechanism design and the revelation principle}
\label{sec:mech:revelation}

A \emph{general mechanism}~\citep{Myerson1981} is a
type-space-and-action-rule
$M = (\Theta, \Act, g : ({\textstyle\prod_i}\Act_i) \to (\Players \to \RR))$;
a strategy profile is a family $(\sigma_i : \Theta_i \to \Act_i)_i$.
A \emph{direct mechanism} is a mechanism with $\Act_i = \Theta_i$;
the canonical strategy is the identity.  The Lean definition
of Bayesian incentive compatibility is the ex-ante finite condition
that truthful reporting is a Bayes-Nash equilibrium, quantified over
type-contingent misreport functions
$s'_i : \Theta_i \to \Theta_i$.

\begin{theorem}[Revelation principle]
\label{thm:revelation}
Let $M$ be a general mechanism with finite types and actions, $\mu$
a common prior over types, and $\sigma$ a Bayesian Nash equilibrium
of $M$.  Then the direct mechanism
$M^{\mathrm{direct}}$ obtained by composing each player's report with
$\sigma_i$ — i.e.\ $M^{\mathrm{direct}}(\theta) = g(\sigma(\theta))$ —
is BIC against type-contingent misreports.
\end{theorem}

\noindent The proof is the textbook argument in finite, ex-ante
form:\ a type-contingent report $s'_i$ in $M^{\mathrm{direct}}$
corresponds to the type-contingent deviation
$\theta_i \mapsto \sigma_i(s'_i(\theta_i))$ in $M$, and the Nash
hypothesis on $\sigma$ rules out profitable deviations.

\leancorr{\TT{Mechanism/MechanismDesign.lean} and
\TT{Mechanism/RevelationPrinciple.lean} (\TT{GeneralMechanism},
\TT{IsBNE}, \TT{toDirect}, \TT{revelation\_principle}).
Supplement~§8.3.}

\paragraph{Vocabulary and surrounding lemmas.}\
The library exposes \TT{isStrategyProof} as a synonym for
\TT{isIC} (the literature uses the two interchangeably).
\TT{Mechanism.socialWelfare} is the utilitarian aggregate of
per-player payoffs on a report profile, and
\TT{IsIndividuallyRational} is the participation constraint
parameterized by an outside-option function $o : \Players \to
\RR$:\ every player gets at least $o_i$ at every type profile.
The standard zero-outside-option normalization is exposed as
\TT{IsNonNegative}.  These are scaffolding for downstream
welfare-vs-truthfulness analyses.

\subsection{Dominant-strategy implementability}
\label{sec:mech:implementability}

The Bayesian mechanism above is payoff-only:\ its rule returns a
payoff vector and never names a chosen alternative.  A second
mechanism layer, \TT{SCFWithPayments}, makes the \emph{allocation}
explicit — agents have type-dependent valuations
$v_i : \Theta_i \to A \to \RR$ over a shared alternative set $A$, a
choice rule selects an alternative, and a payment rule charges each
agent — so that quasilinear utility
$v_i(\theta_i, \mathrm{choice}(\theta)) - p_i(\theta)$ is defined and
implementability acquires structural content.

\paragraph{Necessity:\ weak monotonicity.}\
A choice rule is \emph{weakly monotone} when, holding the others'
reports fixed, an agent's value gain from the alternative chosen at one
of its types over the alternative chosen at another is larger under the
first type.  Every dominant-strategy incentive-compatible mechanism has
a weakly monotone choice rule:\ adding the two truthfulness
constraints — type $\theta_i$ not wishing to report $\theta'_i$, and
$\theta'_i$ not wishing to report $\theta_i$ — cancels the payments and
leaves exactly weak monotonicity.  This is the multi-dimensional
analogue of Myerson's monotonicity~\citep{Bikhchandani2006,
Myerson1981} and the necessary half of the implementability
characterization; the converse (weak monotonicity suffices on convex
domains) is not formalized.

\paragraph{Single-parameter Myerson payments.}\
For real single-parameter direct mechanisms, the library proves the
one-dimensional implementability theorem in the envelope-payment form.
An allocation rule $x_i(\theta)$ is monotone in agent $i$'s own report
when increasing $\theta_i$ weakly increases $x_i$.  The Myerson
payment rule is
\[
  p_i(\theta)
    =
  \theta_i x_i(\theta)
    -
  \int_{0}^{\theta_i}
      x_i(z,\theta_{-i})\,dz .
\]

\begin{theorem}[Myerson payment formula and implementability]
\label{thm:myerson-payments}
Every DSIC single-parameter mechanism has a monotone allocation rule.
Conversely, every monotone allocation rule whose one-dimensional slices
are interval-integrable is implemented in dominant strategies by the
Myerson payment rule.  If the mechanism is zero-normalized and the
allocation slices are continuous, then DSIC payments are uniquely
given by the same envelope formula.
\end{theorem}

\noindent The necessity proof is the standard payment-sandwich
argument: truthfulness at reports $\theta_i$ and $b_i$ gives upper and
lower bounds on the payment difference, and their compatibility forces
monotonicity.  The sufficiency proof integrates the monotone
allocation slice and uses the envelope identity to show that every
misreport loses weakly relative to truth-telling.

\leancorr{\TT{Mechanism/Myerson.lean}
(\TT{SingleParameterMechanism}, \TT{payment\_sandwich},
\TT{isMonotone\_of\_isDSIC},
\TT{withMyersonPayment\_isDSIC\_of\_isMonotone},
\TT{payment\_eq\_myersonPayment\_of\_isDSIC\_of\_zeroNormalized}, and
\TT{existsUnique\_zeroNormalized\_payment\_of\_isMonotone}).}

\paragraph{Sufficiency:\ affine maximizers and VCG.}\
An \emph{affine maximizer} fixes positive weights $w_i$ and an
alternative bias $\kappa : A \to \RR$ and selects the alternative
maximizing $\sum_i w_i\,v_i(\theta_i, \cdot) + \kappa$, paired with
Clarke pivot payments:\ each agent is charged the externality it
imposes on the others — their best attainable weighted welfare minus
their welfare at the chosen alternative — a nonnegative transfer.
Scaling agent $i$'s utility by $w_i$ recovers the entire affine
objective at the chosen alternative — which the choice rule maximizes —
minus a charge independent of $i$'s own report, so no misreport can
help and the mechanism is dominant-strategy truthful.  By Roberts'
theorem~\citep{Roberts1979} affine maximizers are essentially the only
dominant-strategy rules on unrestricted domains; the library proves the
truthful direction, with \emph{VCG} (unit weights, no bias:\ efficient
allocation with Clarke pivot payments) as the canonical special case,
recovering the auction-layer VCG result of \Cref{sec:mech:auctions} at
the abstract level.

\leancorr{\TT{Mechanism/Monotonicity.lean}
(\TT{SCFWithPayments}, \TT{IsDSIC}, \TT{IsWeaklyMonotone},
\TT{weaklyMonotone\_of\_dsic}); \TT{Mechanism/AffineMaximizer.lean}
(\TT{affineMaximizer}, \TT{affineMaximizer\_isDSIC},
\TT{affineMaximizer\_payment\_nonneg}, \TT{vcg}, \TT{vcg\_isDSIC}).}

\subsection{K-implementation:\ influencing a game by nonnegative
transfers}
\label{sec:mech:kimplementation}

Mechanism design creates a game; \emph{k-implementation}
\citep{MondererTennenholtz2004} starts from a game that already
exists.  An interested party — unable to redesign the game, prohibit
strategies, or charge the players — observes the realized strategy
profile and commits to nonnegative transfers
$V : \prod_i X_i \to \RR_{\ge 0}^{\Players}$ on top of the players'
utilities.  Rationality is one-round dominance:\ players avoid
dominated strategies, where throughout this layer ``dominates'' means
weak dominance with a strict witness.  A transfer scheme
\emph{implements} a target set $O$ when the subsidized game has at
least one undominated profile and every undominated profile lies in
$O$; it is a \emph{$k$-implementation} when in addition the total
transfer at every surviving profile is at most $k$; the
\emph{implementation price} of $O$ is the infimum of feasible budgets.
The point of the notion is that promises off the desired path cost
nothing, so desirable outcomes are often implementable far below the
face value of the promises.  The library defines subsidies and
implementation on arbitrary kernel games — no finiteness enters the
definitions — and proves the singleton theory in a strengthened form.

\begin{theorem}[Singleton implementation price]
\label{thm:kimpl-singleton}
For a kernel game with finitely many players and finite nonempty
strategy sets, the implementation price of a singleton target
$\{z\}$ is
\[
  k(z) \;=\; \sum_{i}\; \max_{x_i \in X_i}
      \bigl( u_i(x_i, z_{-i}) - u_i(z) \bigr).
\]
The formula holds as a greatest lower bound with no further
hypotheses; the infimum is attained — and consequently $z$ is Nash iff
$\{z\}$ is 0-implementable — under the nondegeneracy condition that
every player has \emph{another} player with an alternative to its
target action.
\end{theorem}

\noindent Both refinements sharpen the original paper, which imposes
nondegeneracy on the whole statement.  An
$\varepsilon$-bump construction shows the greatest-lower-bound form is
unconditional, and a machine-checked one-player counterexample shows
where nondegeneracy is genuinely load-bearing:\ in a one-player game
with a payoff tie the target is Nash and yet no zero-budget
implementation exists, because without an opponent no strict witness
can dominate the off-target action.  Only the direction
``0-implementable $\Rightarrow$ Nash'' is hypothesis-free.  A second
subtlety justifies the finiteness:\ in infinite games a dominated
strategy need not be dominated by any \emph{undominated} strategy
(the paper's infinite-ascent example, mechanized both as an
order-theoretic core and as a concrete game), so the nonemptiness
clause of implementation can fail catastrophically.  The finite case
is repaired once and for all by a reusable maximal-element lemma for
transitive irreflexive relations on finite types.

\leancorr{\TT{Core/Subsidy.lean}; \TT{Implementation/Basic.lean}
(\TT{IsImplementation}, \TT{IsKImplementation}, \TT{implPrice});
\TT{Implementation/Singleton.lean} (\TT{implementationGap},
\TT{singletonTransfer\_isKImplementation}, \TT{implPrice\_singleton},
\TT{exists\_zero\_singletonKImplementation\_iff\_isNash});
\TT{Implementation/Examples.lean}
(\TT{onePlayerTie\_zero\_not\_mem\_implementationCosts});
\TT{Concepts/Dominance/Undominated.lean}
(\TT{InfiniteAscentCounterexample}); \TT{Math/Fintype.lean}
(\TT{exists\_maximal\_rel\_of\_finite\_trans\_irrefl}).}

\paragraph{A price calculus.}\
Because the singleton price is a closed-form functional of the
deviation gaps, structural corollaries come cheaply.  In an exact
potential game~\citep{MondererShapley1996} the gaps telescope through
the potential $\Phi$:\ the price of $\{z\}$ equals
$\sum_i \max_{x_i} (\Phi(x_i,z_{-i}) - \Phi(z))$, is bounded by
$n\,(\max \Phi - \Phi(z))$, and vanishes exactly at the local maxima
of $\Phi$ under unilateral moves — the Nash-iff-free-implementation
equivalence becomes a discrete first-order condition.  In a different
direction, if the game's outcome kernel is
$\varepsilon$-differentially private — a unilateral strategy change
multiplies every outcome probability by at most $e^{\varepsilon}$ —
and utilities lie in $[0,U]$, then every deviation gain is at most
$(e^{\varepsilon}-1)U$, so \emph{every} singleton target has price at
most $n\,(e^{\varepsilon}-1)\,U$.  This recasts the
approximate-truthfulness argument of
\citet{McSherryTalwar2007} as a uniform implementation-price bound; a
zero-utility game with a fully revealing kernel witnesses that the
converse fails.

\leancorr{\TT{Implementation/Singleton.lean}
(\TT{implPrice\_singleton\_eq\_sum\_potentialGaps},
\TT{implPrice\_singleton\_le\_card\_mul\_potentialGap},
\TT{implPrice\_singleton\_eq\_zero\_iff\_potential\_localMax});
\TT{Implementation/DifferentialPrivacy.lean}
(\TT{IsEpsilonDifferentiallyPrivate},
\TT{implPrice\_singleton\_le\_card\_mul\_dp\_bound});
\TT{Implementation/Examples.lean}
(\TT{revealingZeroUtilityGame\_not\_zeroDifferentiallyPrivate}).}

\paragraph{Mixed extensions, and a corrected Theorem 3.}\
When the interested party can observe mixed strategies themselves —
the paper's ``algorithm observability'' — transfers attach to mixed
profiles, and the paper's Theorem~3 asserts that a mixed profile is a
mixed equilibrium \emph{iff} it is 0-implementable in the mixed
extension, deriving both directions from its (unproved) regular-games
Theorem~2.  The library proves the constructive direction:\ every
observable mixed Nash profile is 0-implementable.  The converse is
\emph{false} as stated:

\begin{theorem}[Failure of the converse for unrestricted transfers]
\label{thm:kimpl-mixed-converse}
There is a two-player $2\times2$ game and a non-Nash profile $z$ of
its mixed extension together with nonnegative transfers on mixed
profiles under which the subsidized mixed extension has exactly one
undominated profile, namely $z$, at zero transfer — a
0-implementation of a non-equilibrium.
\end{theorem}

\noindent The construction rewards proximity to the target uniformly
across opponent strategies, so every non-target mix is strictly
dominated by the mix halfway closer to the target — an infinite
ascent with no maximum below the target — while a separate
off-column bonus protects the target itself and vanishes on the
desired path.  The transfer is necessarily discontinuous at the
target; the diagnosis is that the paper's Theorem~2 (and hence the
converse half of Theorem~3) requires an implicit admissibility class
for transfers.  To our knowledge this correction is new:\ the
published corrections to the paper
\citep{Eidenbenz2007, Eidenbenz2011, DengTangZheng2016} concern its
algorithmic claims only.

The \emph{corrected} converse is then proved under a semantic
admissibility hypothesis — for each player, the subsidized mixed game
must admit an \emph{undominated best response} to the target
opponents, which is exactly the selection property whose finite-game
analogue is the finite-ascent lemma.  Under it, the implementation
price of a mixed singleton equals the sum of pure-deviation gaps (the
affine-in-own-strategy structure collapses mixed deviations to a
finite maximum), the infimum is attained under the mixed
nondegeneracy condition via a distance-bump transfer whose strict
witnesses avoid the counterexample's discontinuous jump, and
zero-cost admissible implementation is equivalent to mixed Nash.  The
counterexample's transfer is proven to violate the selection
hypothesis, so the two results fit together exactly.  A concrete
\emph{analytic} sufficient condition is also formalized, resting on a
game-free compactness lemma of independent interest:\ on a compact
own-strategy space, upper-semicontinuity of every payoff column
implies that every strategy has an undominated weak dominator (Zorn's
lemma plus the finite-intersection property — no continuity in the
column variable is needed).  Consequently columnwise own
upper-semicontinuity of the subsidized payoff on the finite mixed
simplex alone implies the selection property, and the same analytic
price formula holds.  Budgets strictly above the formula are realized
by continuous positive-slack distance-penalty transfers; under the
same mixed nondegeneracy condition as above, exact attainment is by
the distance-bump transfer, whose subsidized own-payoff is a
continuous base function with a single upper-semicontinuous spike at
the target.  The regular-game interface packaging itself is also
formalized; what remains is a Theorem~2-style price statement for
arbitrary compact strategy spaces.

\leancorr{\TT{Math/Topology/WeakDominance.lean}
(\TT{exists\_pointwiseUndominated\_dominator\_of\_compact\_usc},
\TT{UpperSemicontinuous.update\_of\_le});
\TT{Implementation/Mixed.lean}
(\TT{mixedNash\_exists\_zero\_singletonKImplementation},
\TT{MixedSelectionAdmissibleAt}, \TT{mixedAdmissibleImplPrice\_singleton},
\TT{mixedSingletonDistanceTransfer},
\TT{exists\_zero\_mixedAdmissibleKImplementation\_iff\_isNash},
\TT{MixedOwnUpperSemicontinuous},
\TT{MixedOwnUpperSemicontinuous.selectionAdmissible},
\TT{mixedAnalyticImplPrice\_singleton},
\TT{mixedAnalyticImplPrice\_singleton\_of\_nondegenerate},
\TT{exists\_zero\_mixedAnalyticKImplementation\_iff\_isNash},
\TT{mixedSingletonDistanceTransfer\_analyticAdmissible},
\TT{mixedSingletonContinuousApproxTransfer});
\TT{Implementation/Regular.lean}
(\TT{RegularOwnUpperSemicontinuous},
\TT{RegularOwnUpperSemicontinuous.selectionAdmissibleAt});
\TT{Implementation/Examples.lean}
(\TT{mixedConverse\_zeroKImplementation\_and\_not\_nash},
\TT{mixedConverse\_implPrice\_singleton\_eq\_zero},
\TT{mixedConverseTransfer\_not\_selectionAdmissible}).}

\paragraph{Implementation devices and correlated equilibria.}\
Without algorithm observability the paper's Section~6 interposes a
\emph{recommendation device}:\ signals are drawn from a target
distribution $\mu$, each player privately receives its coordinate,
and transfers may depend on the signal profile and the realized
actions.  The library compiles the device and game into a kernel game
whose strategies are signal-contingent rules, relates the paper's
per-signal (interim) dominance to ex-ante dominance of the compiled
game by a law-of-total-expectation bridge, and proves the paper's
Theorems~7 and~8:\ for a correlated equilibrium $\mu$ (in particular,
for the product law of a mixed Nash profile), the obedient profile is
implemented \emph{exactly}, at zero transfer on every supported
signal, with obedience the unique undominated rule profile.
Uniqueness needs full support of $\mu$ — a rule deviating only at a
never-drawn signal is observationally obedient — and an available
non-conforming opponent action; both hypotheses are discharged on a
concrete uniform-correlated example.  The bonus that enforces
obedience is paid only when some opponent deviates, hence never on
path.

Beyond the paper, the layer distinguishes the pointwise realized
budget from the mediator's \emph{expected} outlay and proves a
quantitative generalization to arbitrary target laws:\ any device
whose obedient profile is interim dominant for $\mu$ pays in
expectation at least $\sum_i \max_{d_i}
\operatorname{swapRegret}_i(\mu,d_i)$ — so expected-zero-cost
implementation forces $\mu$ to be a correlated equilibrium — while a
regret-paying device (conditional swap-regret payments on path, the
usual bonus off path) implements every finite $\mu$ at exactly that
expected budget.  Thus the formal expected-cost device theorem has the
closed regret expression \(\sum_i \max_{d_i}
\operatorname{swapRegret}_i(\mu,d_i)\), and the zero-cost correlated
equilibrium theorem is the regret-zero case, connecting the device
layer to the regret characterization of correlated equilibria in
\Cref{sec:solution-concepts}.  Under the paper's \emph{pointwise}
budget the price is instead sandwiched between the largest one-signal
conditional regret and the largest supported sum of conditional
regrets, with equality at the lower end when regret obligations do
not overlap on the support; the general pointwise optimum is a finite
linear program, not a closed regret expression.  The linear program
itself is formalized (budget rows plus conditional-incentive rows over
on-path payments), with reusable weighted dual certificates for price
lower bounds, and a concrete overlapping instance whose exact price
$4/3$ is machine-checked to lie strictly inside the sandwich $[1,2]$.

\leancorr{\TT{Implementation/Device.lean}
(\TT{InterimDominantKImplementsDistribution},
\TT{exists\_recommendationDevice\_obedient\_zeroTransfer\_exactKImplementation},
\TT{InterimWeaklyDominates.toWeaklyDominates});
\TT{Implementation/Examples.lean}
(\TT{deviceTie\_exists\_strictRecommendationDevice\_exactKImplementation});
\TT{Implementation/Device.lean}
(expected-cost side:\ \TT{expectedRealizedTransfer},
\TT{sum\_maxSwapRegret\_le\_expectedRealizedTransfer\_of\_interimDominant},
\TT{isCorrelatedEq\_of\_interimDominantExpectedKImplDist\_zero},
\TT{regretRecommendationDevice\_obedient\_interimDominantExpectedKImplDist};
pointwise side:\ \TT{recommendationDevicePointwiseImplPrice\_sandwich},
\TT{recommendationDevicePointwiseImplPrice\_eq\_recommendationPaymentImplPrice},
\TT{le\_recommendationPaymentImplPrice\_of\_weighted\_certificate});
\TT{Implementation/DeviceLP.lean}
(\TT{recommendationPaymentFeasible\_iff\_lpFeasible},
\TT{le\_recommendationPaymentImplPrice\_of\_dual\_certificates});
\TT{Implementation/Examples.lean}
(\TT{OverlappingRegretPaymentExample.recommendationPaymentImplPrice\_strict\_between\_sandwich});
\TT{Concepts/Correlation/Regret.lean} (\TT{maxSwapRegret},
\TT{maxConditionalSwapRegret}).}

\paragraph{Signal-blind transfers, VCG, and dominance rigidity.}\
For incomplete-information (informational) games the party may see
actions but not the realized signals, so transfers are
\emph{signal-blind}.  Zero-cost implementation is then impossible in
general:\ the paper's Figure~4 game, transcribed payoff-by-payoff, has
an ex-post equilibrium that no signal-blind transfer scheme
$k$-implements for \emph{any} $k$.  The positive result (the paper's
Theorem~6) concerns VCG:\ in a combinatorial auction whose valuations
range over a quasi-field, a frugal surplus-maximizing VCG mechanism
admits a conformance bonus making the bundled truthful profile
ex-post dominant at zero cost — the library proves the underlying
Lemma~1 outright and shows by two further counterexamples that its
hypotheses are necessary (frugality, via a four-good example where a
$\Sigma$-based cheat strictly gains; bounded valuations, via a
scaling argument).  The paper's stronger claim that the truthful
profile is the \emph{unique} undominated one is governed by a general
rigidity fact:

\begin{theorem}[Dominance rigidity]
\label{thm:kimpl-rigidity}
If a strategy weakly dominates every strategy of a player, then that
player's undominated strategies are exactly those with identical
payoff against every opponent context.  Consequently, when each
coordinate of a profile weakly dominates everything, the undominated
profiles form the product of these payoff-equivalence classes, and
the target is uniquely undominated iff every class is trivial.
\end{theorem}

\noindent This turns the paper's uniqueness claim into an explicit
non-redundancy condition on the mechanism — two reports that induce
the same allocation and payments against every opponent report are
necessarily \emph{both} undominated, witnessed on a concrete
redundant-report auction — and replaces bespoke strictness
hypotheses in the device and VCG capstones by an unconditional
characterization of the undominated set.

\leancorr{\TT{Mechanism/InformationalGame.lean}
(\TT{StrategyEquivalentWithTransfer},
\TT{IsExPostDominantProfileWithTransfer.undominatedStrategyProfiles\_eq});
\TT{Implementation/Informational.lean}
(\TT{figure4\_no\_signalBlind\_kImplementation},
\TT{conformanceBonusTransfer\_isSignalBlindKImplementation\_of\_strict},
\TT{conformanceBonusTransfer\_undominatedStrategyProfiles\_eq});
\TT{Implementation/VCG.lean}
(\TT{frugal\_vcg\_bundled\_report\_ge\_qBased\_report},
\TT{conformanceBonusTransfer\_bundledTruthful\_undominatedStrategyProfiles\_eq},
\TT{lemma1\_fails\_without\_frugality\_trueUtility},
\TT{not\_exists\_uniform\_bound\_informationalGame});
\TT{Concepts/Dominance/Undominated.lean} (\TT{UtilityEquivalent},
\TT{undominatedProfiles\_eq\_utilityEquivalentClass\_of\_forall\_weaklyDominates}).}

\paragraph{The corrections literature.}\
Two of the paper's algorithmic claims were corrected later, and the
library mechanizes the corrected picture where it can be
self-contained.  The exact-implementation algorithm~OP is not
optimal~\citep{Eidenbenz2007, Eidenbenz2011}:\ on the published
counterexample the library proves \emph{both} sides — the optimal
exact cost is~$3$, and the OP-shaped transfer family has minimum
cost exactly~$5$ — so the suboptimality of the algorithm's output,
not merely of one transfer, is machine-checked.  The paper's
NP-hardness proof by a two-player SAT reduction is incorrect
\citep{Eidenbenz2007}; the corrected reductions of
\citet{DengTangZheng2016} (six players for pure dominance, two
players for dominance by mixed strategies) are the natural target for
a future \TT{SATReduction} module and are not yet formalized, as they
depend on the external constructions.

\tradeoff{Three observability regimes are deliberately distinguished:\
transfers on realized pure profiles (the base theory), transfers on
mixed profiles (Theorem~3's ``algorithm observability,'' where the
corrected converse above lives), and signal-blind transfers on
actions (devices, Figure~4, VCG).  The paper's Theorem~2 for
compact-continuous ``regular'' games is formalized in the part needed
for the mixed-extension theorem:\ the semantic selection hypothesis,
the compact upper-semicontinuity bridge that implies it, and the
regular-game interface packaging.  What is not formalized is a full
Theorem~2-style price statement for arbitrary compact strategy
spaces.  The mixed-extension counterexample above is exactly the
witness that the paper's omitted proof needs an admissibility class
for transfers; unrestricted transfers make the converse false.}

\subsection{Hidden-action contracts}
\label{sec:mech:contracts}

The principal–agent layer \TT{Mechanism/Contracts/Basic.lean}
formalizes the moral-hazard problem~\citep{Holmstrom1979}:\ a
risk-neutral agent privately chooses a costly action that
stochastically determines an outcome, and the principal designs a
payment schedule on outcomes.  The library records the agent's expected
utility (expected payment minus action cost), the principal's (expected
outcome reward minus expected payment), the welfare identity that the
two sum to social surplus
(\TT{principalUtility\_add\_agentUtility}), limited-liability and
incentive predicates, and \emph{linear contracts}
$t \mapsto \alpha\,r(t)$ with the principal-share formula
$(1-\alpha)\,\mathbb{E}[\mathrm{reward}]$
(\TT{principalUtility\_linearPayment}).  The \emph{first-best} benchmark
is the headline:\ moral hazard cannot beat the full-information surplus
(\TT{principalUtility\_le\_firstBest}).  The ``sell the firm'' contract
($\alpha = 1$) makes the agent the residual claimant, so its privately
optimal action is a surplus-maximizer
(\TT{isIncentivized\_linearPayment\_one\_iff}) — first-best
\emph{efficiency} — though the principal then extracts none of it, its
payoff falling to zero (\TT{principalUtility\_linearPayment\_one}).

\leancorr{\TT{Mechanism/Contracts/Basic.lean}
(\TT{PrincipalAgent}, \TT{agentUtility}, \TT{principalUtility},
\TT{linearPayment}, \TT{principalUtility\_add\_agentUtility},
\TT{principalUtility\_le\_firstBest}).}

\subsection{Social choice}
\label{sec:mech:social}

The library formalizes the basic social-choice setup:\ a finite
set of alternatives, a finite set of voters with strict
preferences, and a social welfare function or social choice
function, together with the structural fragments needed downstream
(preference profiles, the existence of a social-welfare-maximizing
alternative under utilitarian aggregation).  On top of this it
proves \emph{Arrow's impossibility theorem}~\citep{Arrow1951}:\
with at least three alternatives and a nonempty finite electorate,
every collectively rational social-welfare function satisfying weak
Pareto and independence of irrelevant alternatives (on the domain of
strict rankings) is dictatorial.  The proof follows Geanakoplos's
pivotal-voter argument, and a companion statement records the
dictator's strict preference as coinciding with society's on every
ranking profile.  The strategic analogue, the
Gibbard–Satterthwaite theorem~\citep{Gibbard1973, Satterthwaite1975}
— a strategy-proof social choice function with at least three
alternatives in its range and a nonempty finite electorate is
dictatorial — is also formalized, by the classical reduction to
Arrow through an induced social-welfare function.

The library also formalizes the classical results around majority
rule.  \emph{May's theorem}~\citep{May1952} characterizes the rule
axiomatically:\ over two alternatives, and for an arbitrary finite
electorate in which voters may be indifferent, a social decision
function agrees with majority rule --- the sign of the vote tally ---
on every profile if and only if it is anonymous, neutral, and
positively responsive (\TT{may\_theorem}).  \emph{Condorcet's paradox}
--- pairwise majority preference can
cycle even when every voter has a strict ranking --- is recorded as
the acyclicity failure of the majority relation
(\TT{condorcetMajority\_not\_acyclic}).  Its positive counterpart is
\emph{Black's median-voter theorem}~\citep{Black1948}:\ when
preferences are single-peaked, the median voter's ideal point is a
Condorcet winner (\TT{median\_is\_condorcet\_winner}).  \emph{Sen's
liberal paradox}~\citep{Sen1970} completes the impossibility group:\
no social decision function is simultaneously Paretian and minimally
liberal --- each of two individuals decisive over some pair of
alternatives (\TT{sen\_paretian\_liberal}).

\leancorr{\TT{Mechanism/SocialChoice.lean}, \TT{Mechanism/Arrow.lean},
\TT{Mechanism/GibbardSatterthwaite.lean}, \TT{Mechanism/Condorcet.lean},
\TT{Mechanism/MedianVoter.lean}, \TT{Mechanism/May.lean},
\TT{Mechanism/SenParetianLiberal.lean}.}

\subsection{Fair division}
\label{sec:mech:fair-division}

Fair division is organized under the mechanism/social-choice branch
because its objects are allocations and fairness predicates rather
than strategic equilibria.  The library has separate indivisible and
divisible packages.

\paragraph{Indivisible goods.}\
For finite additive valuations over indivisible goods, the library
defines envy-freeness, EF1, EFX, proportionality, and maximin-share
guarantees.  It proves the standard implications
envy-free $\Rightarrow$ EFX $\Rightarrow$ EF1 under nonnegative
valuations, and envy-free $\Rightarrow$ proportionality when all goods
are allocated.

\begin{theorem}[EF1 for indivisible goods]
\label{thm:indivisible-ef1}
For any finite set of agents and finite set of indivisible goods with
additive valuations, the envy-cycle allocation rule and the
round-robin allocation rule both return allocations that are EF1.
\end{theorem}

\begin{theorem}[Two-agent EFX]
\label{thm:indivisible-efx-two}
For two agents with additive nonnegative valuations over finitely many
indivisible goods, there exists an EFX allocation.
\end{theorem}

\begin{theorem}[MMS upper bound]
\label{thm:mms-average}
For additive nonnegative valuations, each agent's maximin share is at
most its proportional average value for the whole set of goods.
Consequently every proportional allocation is a maximin-share
allocation.
\end{theorem}

\leancorr{\TT{Mechanism/FairDivision/Indivisible/Basic.lean}
(\TT{IsEnvyFree}, \TT{IsEF1}, \TT{IsEFX},
\TT{IsEnvyFree.isProportional}, \TT{exists\_efx\_two\_agents});
\TT{Indivisible/EnvyCycle.lean}
(\TT{envyCycleAllocation\_isEF1});
\TT{Indivisible/RoundRobin.lean}
(\TT{roundRobinAllocation\_isEF1});
\TT{Indivisible/MaximinShare.lean}
(\TT{maximinShare\_le\_average}, \TT{IsProportional.isMaxminShare}).}

\paragraph{Divisible goods.}\
The divisible package treats cake-cutting on the unit interval with
finite non-atomic measure valuations.  The two-agent theorem is
cut-and-choose in the Steinhaus tradition~\citep{Steinhaus1948};
proportionality for any finite number of agents is proved by a
Dubins--Spanier-style moving-knife construction
\citep{DubinsSpanier1961}; envy-free existence for any finite number
of agents is proved through Stromquist's KKM/compactness route
\citep{Stromquist1980}.

\begin{theorem}[Divisible fair division]
\label{thm:divisible-fair-division}
Let $N$ be a nonempty finite set of agents, and let each agent have a
finite non-atomic measure on the unit interval.  Then there exists a
measurable allocation of the interval that is envy-free.  There also
exists a proportional allocation, and every envy-free allocation is
proportional under the library's finite-agent proportionality
predicate.
\end{theorem}

\leancorr{\TT{Mechanism/FairDivision/Divisible/CutAndChoose.lean}
(\TT{cutAndChoose\_ef\_exists}, \TT{cutAndChooseRule\_isEnvyFree});
\TT{Divisible/DubinsSpanier.lean}
(\TT{dubinsSpanierProportional},
\TT{dubinsSpanierRule\_isProportional});
\TT{Divisible/Stromquist.lean}
(\TT{stromquist\_envyFree\_exists});
\TT{Divisible/Existence.lean}
(\TT{proportional\_exists}, \TT{ef\_exists},
\TT{ef\_and\_proportional\_exists}).}

\subsection{Auctions}
\label{sec:mech:auctions}

The auction directory contains six formats at different
levels of completeness:

\begin{itemize}\itemsep4pt
\item \emph{Vickrey (second-price)}~\citep{Vickrey1961}.  The
highest bidder wins and pays the second-highest bid; truthful
bidding is a weakly dominant strategy.  The proof is by case
analysis on whether $i$'s valuation exceeds the maximum other bid; in
each case, the EU of bidding $v_i$ matches or exceeds the EU of
any other bid.

\item \emph{First-price sealed bid}~\citep{Vickrey1961}.  The
highest bidder pays their own bid.  The file defines the
deterministic first-price game and proves that no single bid is a
dominant strategy for any bidder.

\item \emph{All-pay auction}~\citep{KrishnaMorgan1997}.  Every
bidder pays their own bid regardless of winning.  The
formalization does not build a kernel game; it records elementary
payoff facts such as overbidding being unprofitable, winner
profit, rent dissipation, and a symmetric negative-payoff lemma.

\item \emph{Reserve Vickrey}.  A second-price auction with a reserve
price is packaged as a kernel game via the principled auction-game
interface.  The library proves dominant-strategy truthfulness and
records the usual utility identities for winners, losers, and
allocation.

\item \emph{Vickrey–Clarke–Groves (VCG)}~\citep{Vickrey1961, Clarke1971,
Groves1973}.  Generalizes Vickrey to multi-item, combinatorial
settings:\ each agent's payment equals the externality they impose
on others.  The library proves dominant-strategy incentive
compatibility for finitely many players, with the efficient
allocation supplied as a hypothesis and no finiteness imposed on the
allocation or type spaces.

\item \emph{Knapsack auctions}.  Agents have public weights and
single-parameter values; feasible binary allocations are subsets whose
total weight is at most capacity.  The welfare-maximizing allocation
rule is monotone in each agent's reported value, so the Myerson
payment rule makes the mechanism DSIC.  The same file also includes a
finite dynamic-programming optimal allocation for natural weights and
values, and a half-approximation theorem for the integral greedy rule
against the dynamic-programming optimum.
\end{itemize}

Vickrey and first-price are direct kernel-game constructions via
the \TT{ofEU} adapter.  VCG is packaged as a finite setup with an
efficient allocation rule and a payment function; all-pay is
packaged as a collection of real-valued payoff lemmas.  Solution
concepts come from \Cref{sec:solution-concepts} wherever a full
kernel game has been constructed.

\begin{theorem}[Knapsack DSIC and approximation]
\label{thm:knapsack-auction}
For a nonnegative-capacity binary knapsack auction, the
welfare-maximizing allocation rule is monotone and the mechanism using
Myerson payments is DSIC.  For natural weights and values, the dynamic
programming allocation is feasible and optimal, and the integral
greedy value is at least one half of the dynamic-programming optimum.
\end{theorem}

\tradeoff{Auctions are typically formalized in the literature
under continuous bid spaces and continuous valuations; this gives
the cleanest expressions for symmetric equilibrium bids and
expected revenue formulas.  Our discrete formalization is
faithful to the strategic content it formalizes
(truthful bidding is dominant in Vickrey/VCG; first-price has no
dominant bid; all-pay payoff inequalities are available) but is
not the primary place a researcher would go to study revenue
ranking under continuous distributions.  We treat the auction layer
primarily as a smoke test of the core library; downstream, an
extension to continuous auctions over mathlib's measure theory
would be the natural next step.}

\leancorr{\TT{Auctions/Vickrey.lean} (\TT{vickrey\_truthful\_dominant}),
\TT{Auctions/ReserveVickrey.lean}, \TT{Auctions/FirstPrice.lean},
\TT{Auctions/AllPay.lean}, \TT{Auctions/VCG.lean}, and
\TT{Auctions/Knapsack/Basic.lean}
(\TT{welfareMaximizingAllocationRule\_isMonotone},
\TT{welfareMaximizingMechanism\_isDSIC},
\TT{dynamicProgrammingOptimalAllocation\_feasible},
\TT{dynamicProgrammingOptimalAllocation\_optimal}, and
\TT{integralGreedy\_halfApprox\_dpOptimal}).  Supplement~§8.4.}
